% !TeX root = main.tex
\thispagestyle{empty}


\docAuthors, Введение в параллельные вычисления.
--- СПб: \docUniveristy, \docYear.
---~\pageref*{LastPage}~с.
\todo[noinline]{FIX}

\medskip
Рецензент: Зилинберг А.Ю., к.т.н., доцент кафедры радиотехнических систем, ФГАОУ ВО \enquote{Санкт-Петербургский государственный университет аэрокосмического приборостроения}.

\medskip
В пособии излагаются основные понятия и определения теории параллельных вычислений.
Рассматриваются основные принципы построения программ на языке \enquote{Си} для многоядерных и многопроцессорных вычислительных комплексов с общей памятью.
Предлагается набор заданий для проведения лабораторных и практических занятий.

\medskip
Учебное пособие предназначено для студентов, обучающихся по магистерским программам направления \enquote{09.03.04 --- Программная инженерия}, \enquote{09.03.01 --- Информатика и вычислительная техника}, и может быть использовано выпускниками (бакалаврами и магистрантами) при написании выпускных квалификационных работ, связанных с проектированием и исследованием многоядерных и многопроцессорных вычислительных комплексов.

\begin{flushright}
    \includegraphics[scale=0.15]{logo/itmo-logo-black-2026}
\end{flushright}

Университет ИТМО --- ведущий вуз России в области информационных и фотонных технологий, один из немногих российских вузов, получивших в 2009 году статус национального исследовательского университета.
В 2013--2020 годах Университет ИТМО --- участник программы повышения конкурентоспособности российских университетов среди ведущих мировых научно-образовательных центров, известной как проект \enquote{5 в 100}.
Цель Университета ИТМО --- становление исследовательского университета мирового уровня, предпринимательского по типу, ориентированного на интернационализацию всех направлений деятельности.

\vspace*{\fill}

\begin{flushright}
    \copyright~\docUniveristy, \the\year%

    \copyright~\docAuthors, \the\year%
\end{flushright}
